\begin{tikzpicture}[
    node distance=0.8cm and 0.5cm,
    box/.style={rectangle, draw=blue!60, fill=white, rounded corners=4pt, minimum width=2.6cm, minimum height=1cm, align=center, font=\sffamily\footnotesize, thick},
    optbox/.style={rectangle, draw=orange!80, fill=orange!5, rounded corners=4pt, minimum width=2.6cm, minimum height=1cm, align=center, font=\sffamily\footnotesize\bfseries, thick},
    modbox/.style={rectangle, draw=gray!60, fill=white, rounded corners=6pt, minimum width=2.2cm, minimum height=0.8cm, align=center, font=\sffamily\footnotesize, thick},
    title/.style={font=\sffamily\large\bfseries, text=blue!80!black},
    arr/.style={-{Stealth[scale=1.2]}, thick, draw=blue!60},
    optarr/.style={-{Stealth[scale=1.2]}, thick, draw=orange!80},
    modarr/.style={<->, dashed, thick, draw=gray!40},
    panel/.style={fill=gray!5, rounded corners=12pt, inner sep=15pt}
]

% Hilfskoordinaten für exakt bündige und gleich große Hintergründe (vertikal)
\node[inner sep=0pt] (TopN) at (1.5, 0.5) {};
\node[inner sep=0pt] (BotN) at (1.5, -10.2) {};
\node[inner sep=0pt] (TopA) at (8.5, 0.5) {};
\node[inner sep=0pt] (BotA) at (8.5, -10.2) {};
% Um 0.2cm nach rechts verschoben für identischen Zwischenraum
\node[inner sep=0pt] (TopM) at (16.8, 0.5) {};
\node[inner sep=0pt] (BotM) at (16.8, -10.2) {};

% ================= NAIVE RAG =================
\node[title] (TitleNaive) at (1.5, 0) {Naive RAG};

\node[box] (N_User) at (0, -1.5) {Nutzer-\\anfrage};
\node[box] (N_Doc) at (3.0, -1.5) {Daten-\\quellen};

\node[box] (N_Idx) at (3.0, -3) {Indexierung};
% Breite 5.6cm und x=1.5 zentriert führt zu perfekten Abschlüssen (links 0 und rechts 3.0)
\node[box, minimum width=5.6cm] (N_Ret) at (1.5, -5) {Retrieval};
\node[box, minimum width=5.6cm] (N_Gen) at (1.5, -7.5) {Generierung (LLM)};
\node[box, minimum width=5.6cm] (N_Out) at (1.5, -9) {Antwort};

\draw[arr] (N_Doc) -- (N_Idx);
\draw[arr] (N_Idx) -- (N_Ret.north -| N_Idx.south);
\draw[arr] (N_User) -- (N_Ret.north -| N_User.south);
\draw[arr] (N_Ret) -- (N_Gen);
\draw[arr] (N_Gen) -- (N_Out);

\begin{scope}[on background layer]
    \node[panel, fit=(TopN) (BotN) (N_User) (N_Doc)] (PanelNaive) {};
\end{scope}


% ================= ADVANCED RAG =================
\node[title] (TitleAdv) at (8.5, 0) {Advanced RAG};

\node[box] (A_User) at (7.0, -1.5) {Nutzer-\\anfrage};
\node[box] (A_Doc) at (10.0, -1.5) {Daten-\\quellen};

\node[optbox] (A_Pre) at (7.0, -3.2) {Pre-Retrieval\\[1mm]\normalfont\scriptsize(z.B. Rewrite, Expand)};
\node[box] (A_Idx) at (10.0, -3.2) {Indexierung};

% Breite 5.6cm und x=8.5 sorgt auch hier für exakte Rechts- und Linksbündigkeit
\node[box, minimum width=5.6cm] (A_Ret) at (8.5, -5) {Retrieval};

\node[optbox, minimum width=5.6cm] (A_Post) at (8.5, -6.5) {Post-Retrieval\\[1mm]\normalfont\scriptsize(Reranking, Kompression)};

\node[box, minimum width=5.6cm] (A_Gen) at (8.5, -8) {Generierung (LLM)};
\node[box, minimum width=5.6cm] (A_Out) at (8.5, -9.5) {Antwort};

\draw[arr] (A_Doc) -- (A_Idx);
\draw[arr] (A_Idx) -- (A_Ret.north -| A_Idx.south);
\draw[optarr] (A_User) -- (A_Pre);
\draw[optarr] (A_Pre) -- (A_Ret.north -| A_Pre.south);
\draw[optarr] (A_Ret) -- (A_Post);
\draw[optarr] (A_Post) -- (A_Gen);
\draw[arr] (A_Gen) -- (A_Out);

\begin{scope}[on background layer]
    \node[panel, fit=(TopA) (BotA) (A_User) (A_Doc)] (PanelAdv) {};
\end{scope}


% ================= MODULAR RAG =================
% Alle X-Koordinaten um +0.2 verschoben
\node[title] (TitleMod) at (16.8, 0) {Modular RAG};

\node[rectangle, draw=blue!30, fill=white, dashed, rounded corners=8pt, minimum width=8.2cm, minimum height=8.2cm, thick] (ModBox) at (16.8, -5.5) {};
\node[anchor=south, font=\sffamily\footnotesize\color{blue!70!black}\bfseries] at (ModBox.north) {Baukastensystem (freie Kombination)};

\node[font=\sffamily\LARGE\bfseries, text=blue!15] at (16.8, -5.5) {RAG};

% Inner ring
\node[modbox, draw=orange!80, fill=orange!5] (M_Ret) at (16.8, -4) {Retrieve};
\node[modbox, draw=orange!80, fill=orange!5] (M_Read) at (16.8, -7) {Read};
\node[modbox] (M_Rewrite) at (14.8, -5.5) {Rewrite};
\node[modbox] (M_Rerank) at (18.8, -5.5) {Rerank};

% Outer ring
\node[modbox] (M_Search) at (16.8, -2.3) {Search};
\node[modbox] (M_Mem) at (16.8, -8.7) {Memory};
\node[modbox] (M_Route) at (14.3, -3) {Routing};
\node[modbox] (M_Predict) at (19.3, -3) {Predict};
\node[modbox] (M_Demo) at (14.3, -8) {Demonstrate};
\node[modbox] (M_Fusion) at (19.3, -8) {Fusion};

% Modulare Verbindungen (Netzwerk-Charakter)
\draw[modarr] (M_Rewrite) edge (M_Ret);
\draw[modarr] (M_Ret) edge (M_Rerank);
\draw[modarr] (M_Rerank) edge (M_Read);
\draw[modarr] (M_Read) edge (M_Rewrite);
\draw[modarr] (M_Ret) edge (M_Read);

\draw[modarr, draw=blue!30] (M_Search) edge (M_Ret);
\draw[modarr, draw=blue!30] (M_Route) edge (M_Rewrite);
\draw[modarr, draw=blue!30] (M_Predict) edge (M_Rerank);
\draw[modarr, draw=blue!30] (M_Mem) edge (M_Read);
\draw[modarr, draw=blue!30] (M_Fusion) edge (M_Rerank);
\draw[modarr, draw=blue!30] (M_Demo) edge (M_Rewrite);

\begin{scope}[on background layer]
    \node[panel, fit=(TopM) (BotM) (ModBox)] (PanelMod) {};
\end{scope}

\end{tikzpicture}